
% 来源：EPUB part221.xhtml
\section{希望自己和后代生活在更好的环境里}

我做公益但不喜欢做慈善这个说法，甚至不觉得自己是慈善的。个人觉得公益和慈善是有区别的。公益感觉自己身在其中，慈善则有点居高临下的感觉。我有做公益的需求，希望自己和后代可以生活在更好的环境里。（2020-08-12）

% 来源：EPUB part222.xhtml
\subsection{公益是每个人自己该做的事情}

\BookQuestion{网友：投资是为了赚钱吗？}

呵呵，这个问题有意思。我想，对大多数人而言，投资当然是为了赚钱吧。投资的过程也是很有乐趣的。

所以对我个人而言，享受过程可能更重要，因为结果已经不那么重要了。不过，好的结果还能帮到别人，也算是个不错的副产品。（2010-04-24）

其实财富过了够用的线之后，性质会发生变化，多出来的都是责任。（2021-01-21）

\BookQuestion{网友：财富达到哪个级别时是你最开心的阶段？}

好吧，这也许是个有意思的问题。记得当年拿到高考分数的时候好像是非常开心的。后来也有过很多开心的时刻，但和赚多少钱似乎都没太大关系。很多年前我就大概猜到我会赚很多钱的，所以赚多少钱我都不会太意外，也没啥惊喜。就像苹果涨了一样。（2024-06-26）

\BookQuestion{网友：如何提高自己的舍得能力，特别是在金钱方面。}

理性，想长远！（2020-12-02）

\BookQuestion{网友：段总一切大小的布施善行值得赞叹。}

呵呵，拿自己用不着的东西去帮帮人真不算什么。（2010-03-09）

保持受捐者的尊严和帮助受捐者一样重要！

虽然人们总是把慈善公益放一起说，但我总觉得公益和慈善是有点区别的，但还没想清楚本质区别是什么。

国内公益环境不太好，公益意识也差很多，相关的法律环境也不太匹配。比如美国的基金会每年的要求是捐5\%，而中国是8\%。 5\%是可持续的概念， 8\%是不太可能持续的。很多年前我问过政策制定者，为什么是8\%，回答是不知道应该是多少，就定了一个比美国高的比例。用几十年和几千几万个公益组织的角度看就明白这个差别有多大了。

我觉得在国内一直坚持做公益的人们真是很伟大，他们所面临的困难比在美国大很多。

顺便说一点公益意识。整体而言，老美的公益意识是很强的，强在大家都觉得公益是每个人自己该做的事情。国内的公益意识总体感觉还是觉得这是富人该干的事，而且由于富人都是“为富不仁”的，所以是赎罪性的，所以不用感谢他们。好消息是，我已经看到国内近几年的巨大进步了，虽然依然还差很远。（2012-05-12）

\BookQuestion{网友：你2008年在浙大设的心平自立贷学金项目现在情况如何？看了那个视频后，我觉得学生们很多没有真正理解你的用意。}

为什么你认为很多学生没有理解呢？我不喜欢帮助“贫困学生”的概念，因为那必须把每个申请人放在放大镜下看他们是不是“足够穷”。

记得在人大时还看到贫困生和捐助人面对面接受捐赠的场景，感觉不太好。不能让人为了得到帮助而损失自尊的。自立贷学金解决了这个问题，由申请人自己决定是不是要借这个钱来帮助自己上学。

这个钱大部分人本来就不需要，但需要的人自然会明白的。

我倒是很希望我上学时可以有这个项目，那样就不用增加父母的负担了。目前好像浙大的留学贷学金需求还蛮大的，第一期的钱已经用完了，第二期的钱应该已经给过去了。人大的情况还没看到。（2011-03-23）

\BookQuestion{网友：王阳明说致良知就是要知善行善，知恶去恶。}

知善行善，知恶去恶。这个说法很好。（2013-03-20）

\BookQuestion{网友：今天刚看到段学长给母校捐款，希望以后也可以向前辈看齐。}

在你触手可及的地方做点公益其实不需要等以后的。（2022-05-16）

\BookQuestion{网友：公益力所能及。只要在做，都是值得称道的。}

是的，哪怕是顺手捡个垃圾。（2022-05-17）

\BookQuestion{网友：我一直有个愿望，将来有钱后帮助那些因家庭穷而辍学的孩子。}

他们辍学一般而言往往不是经济问题，是意识问题，是家长对教育的认识问题，不然就不会有女孩辍学多过男孩的现象。这个问题往往会是恶性循环，一代一代传下去。（2011-02-11）

\BookQuestion{网友：有人会把你的善心当作弱点来利用；有人赌博亏了很多钱负债累累还继续在赌，死不悔改。假如这种人是你的兄弟或亲戚，你还会帮他吗？}

我觉得帮人并不是简单给钱了事，那样会没完没了而且往往没什么好结果，有时可能还会被抱怨给得不够多。所以我们基金会的钱主要是投在教育上。

我帮亲戚也只是帮在教育上。老家的任何直系亲属如果小孩上任何学，家里有困难的我都帮着付学费（有个很好的律师每年都要专门去核实一次）。别的我管不了（除非父母有需求）。（2010-05-14）

\BookQuestion{网友：请问您如何看待教育与家庭支出比例问题？像步步高实验小学这种师资队伍是每个家长都羡慕的。可是作为打工一族未来具有不确定性，如果教育比例支出过大又面临失业的话，这个支出将成为沉重的家庭负担，也害怕因为家庭环境转学影响孩子心理与学习。}

好的教育是昂贵的，更贵的是不好的教育。不过，我们会继续想办法降低员工子弟的入学成本的。（2024-06-21）

% 来源：EPUB part223.xhtml
\subsection{受教育是自助的最好办法}

\BookQuestion{网友：能说说你的慈善心路吗？}

喜欢投资，不会花钱，不想留太多钱给小孩，所以必须找到办法把钱有意义地花出去。花这种钱比赚钱难多了。（2011-04-25）

\BookQuestion{网友：似乎大道在20年前就写好了遗嘱？}

是啊，人有旦夕祸福，早点做准备比较负责任……遗嘱早就写好了，每几年审阅一下。（2024-07-12）

\BookQuestion{网友：有种说法叫养懒汉，确确实实有部分人爱穷，喜欢穷，从一个懒字可以体现出来。你怎样看？}

穷和懒没有必然联系，我见过很多穷人是非常勤快的。富的人也未必就一定勤快，其实我也算个懒人。一个人是否能有财富，取决于很多因素，到底有哪些我也说不清，但受教育程度（不一定是学历啊）肯定是决定因素之一。

这大概就是巴菲特讲的命吧，所以富人应该把财富回馈给社会。富人在赚钱的过程当中已经享受到了很多人生乐趣，多余的财富一定要早早筹划，要找到能高效回馈社会的方法。富人把多余的财富回馈给社会实际上往往要比他们赚到这些财富要难得多得多，再过几十年大家就会在大陆看到很多类似香港小甜甜这样的故事了。（2011-02-08）

济贫的事确实难度比较高，主要是无法鉴别，伦理上也有纠结。所以我们的基金会主要是投教育，觉得授人以渔比较好。上天只助自助之人嘛，受教育是自助的最好办法。顺便说下，我自己确实算懒人，所以会想出或学会许多很好的懒办法。就投资而言，少犯错误其实就是个懒办法。日内交易员们比我可勤快多了。

话说回来，我其实不知道是不是真有懒人。我觉得懒人多数可能是没找到自己喜欢干的事情吧。（2011-02-10）

\BookQuestion{网友：段总在2010年度慈善家捐赠排行榜前十。}

慈善不是比赛，出这种榜的人缺德。（2011-04-28）

\BookQuestion{“企业多做公益慈善事业不吃亏。”（新闻）}

感觉标题把因果关系颠倒了。实际情况应该是这类好企业注重消费者感受，所以能得到消费者好感，从而也就有能力去从事一些公益。其实好企业本身的经营行为从广义的角度看也可以算是公益行为，因为他们在改善人们的生活品质。（2011-12-15）

\BookQuestion{网友：说说我对投资与慈善的看法。1. 慈善也与投资一样要追求“不被关注”，那是最需要的地方。2. 先投资再做慈善或边投资边做慈善，是个人的事没有区别，无需比较，没有高下。都值得尊敬。3. 做慈善是把钱花在最需要的地方，这一点与投资是一样的，投资讲效益，做慈善讲效果。}

1. 任何形式的公益我都不反对。

2. 同意。

3. 公益的效果是投入人心目当中的效果，不是旁观者心目当中的效果。每个人都会希望达到最佳效果，但不能苛求。任何旁观者对效果说三道四都是不道德的，因为长远来讲这会阻止人们做公益。当然公益人士内部的讨论却是有必要的，因为这样可以提高效率。

长远来讲，要求慈善把钱花在最需要的地方的观念是非常不成熟和极端错误的想法，由于没人知道如何才能做到这一点，所以大部分人会选择多一事不如少一事。谁愿意花钱买批评啊？所以我特别希望喜欢对公益慈善说三道四的旁观者闭嘴！（呵呵，不是针对你啊。）刚刚看到马云说要拿出0.3\%的营业额投入到公益当中，作为间接股东，我非常赞成。（2010-05-15）

评价别的公益组织实际上是阻碍人们做公益的一个很重要的原因。你不理解这点，说明你实际没思考过。（2012-07-14）

\BookQuestion{网友：挺多人不了解捐赠，总觉得是为了避税。}

是啊，本金都捐了，避哪门子税？！（2024-07-04）

\BookQuestion{网友：国内做公益真的好难，做好事都怕伤着。}

无欲则刚。（2012-11-05）

国内做公益不容易，官本位的味道很浓，多数人都觉得你做公益肯定是有所图的，经常会觉得有点别扭。（2020-10-09）

\BookQuestion{网友：您知道有哪些比较靠谱的公益途径吗？}

不敢推荐啊，很难了解很深，比了解一个公司都难。不过，即便如此，我还是会继续投钱到公益里面的。一般是先投一点点，然后慢慢增加，如果自己觉得他们做的事情是相对满意的话。当然，本来就很了解的地方容易决定一些，比如母校啥的，但那也不是绝对保险的。做公益是自己的需求，有时候感觉像是做风投。做投资反而要少担心很多。（2024-08-15）
